---
format: yohaku/0.2
target: text
macro: '[[macros/raw]]'
imported: {"bytes":5306,"path":"/root/projects/yohaku/release-prep/yohaku/agent/yohaku_launcher.py","sha256":"443d4a47657fa71f3e4b1c56d7a760f40fd5c39a06b05ee25c43175a3f95b361"}
---

## 正文
```text
#!/usr/bin/env python3
"""Run the existing Yohaku launcher with bounded display-disconnect handling.

Usage in WSL: python3 yohaku-display-check.py --project ~/projects/yohaku
This diagnostic wrapper does not install files or change graphics/sandbox settings.
"""
import argparse
import os
from pathlib import Path
import selectors
import signal
import subprocess
import sys
import time

DISPLAY_ERRORS = {
    b"Io error: Connection reset by peer (os error 104)",
    b"Io error: Broken pipe (os error 32)",
}


def supervise(command, project, output, env=None):
    child = subprocess.Popen(
        command, cwd=project, stdout=subprocess.PIPE,
        stderr=subprocess.STDOUT, start_new_session=True, env=env,
    )
    interrupted = None
    old_handlers = {}

    def interrupt(signum, _frame):
        nonlocal interrupted
        interrupted = signum

    def terminate_group():
        try:
            os.killpg(child.pid, signal.SIGTERM)
        except ProcessLookupError:
            pass
        try:
            child.wait(timeout=1)
        except subprocess.TimeoutExpired:
            pass
        # Native workers can outlive cargo; terminate the owned group as well.
        try:
            os.killpg(child.pid, signal.SIGKILL)
        except ProcessLookupError:
            pass
        child.wait()

    selector = selectors.DefaultSelector()
    consecutive = 0
    pending = b""
    disconnected = False
    exited_at = None

    def line(data):
        nonlocal consecutive, disconnected
        output.write(data)
        output.flush()
        if data.strip() in DISPLAY_ERRORS:
            consecutive += 1
        else:
            consecutive = 0
        disconnected = consecutive >= 4

    try:
        for signum in (signal.SIGINT, signal.SIGTERM):
            old_handlers[signum] = signal.signal(signum, interrupt)
        selector.register(child.stdout, selectors.EVENT_READ)
        while not disconnected and interrupted is None:
            events = selector.select(timeout=0.1)
            if child.poll() is not None:
                if exited_at is None:
                    exited_at = time.monotonic()
                if not events and time.monotonic() - exited_at > 1:
                    break
            if not events:
                continue
            data = os.read(child.stdout.fileno(), 8192)
            if not data:
                break
            pending += data
            while b"\n" in pending and not disconnected:
                current, pending = pending.split(b"\n", 1)
                line(current + b"\n")
            # Even output without newlines must not consume unbounded memory.
            if len(pending) > 65536 and not disconnected:
                line(pending)
                pending = b""
        if pending and not disconnected and interrupted is None:
            line(pending)
        if disconnected:
            output.write(
                b"\nYohaku: repeated display disconnect; stopping this launch.\n"
                b"This does not identify the cause or repair WSLg.\n"
                b"Next diagnostic (in WSL): tail -n 60 /mnt/wslg/weston.log\n"
            )
            output.flush()
            return 70
        if interrupted is not None:
            return 128 + interrupted
        code = child.poll()
        return code if code is not None and code >= 0 else (128 - code if code is not None else 1)
    finally:
        terminate_group()
        selector.close()
        child.stdout.close()
        for signum, handler in old_handlers.items():
            signal.signal(signum, handler)


def launch_environment(backend):
    env = os.environ.copy()
    if backend == "auto":
        try: wsl = "microsoft" in Path("/proc/version").read_text().lower()
        except OSError: wsl = False
        backend = "x11" if wsl else "native"
    if backend == "x11":
        if not env.get("DISPLAY"):
            raise ValueError("X11 display is not available (DISPLAY is unset).")

        # Forward the Fcitx5 input-method environment to X11 applications.
        env["XMODIFIERS"] = "@im=fcitx"
        env["GTK_IM_MODULE"] = "fcitx"
        env["QT_IM_MODULE"] = "fcitx"

        env.pop("WAYLAND_DISPLAY", None)
        env.pop("WAYLAND_SOCKET", None)
    return env, backend


def main():
    parser = argparse.ArgumentParser(description="Launch Yohaku with display-disconnect handling")
    parser.add_argument("--project", type=Path, default=Path(__file__).resolve().parents[1])
    parser.add_argument("--backend", choices=["auto", "native", "x11"], default=os.environ.get("YOHAKU_DISPLAY_BACKEND", "auto"))
    args = parser.parse_args()
    if sys.platform != "linux": parser.error("Run inside WSL/Linux.")
    project = args.project.expanduser().resolve()
    if not (project / "Cargo.toml").is_file(): parser.error(f"Yohaku project not found: {project}")
    try:
        env, backend = launch_environment(args.backend)
        print(f"Yohaku host display: {backend}", flush=True)
        return supervise(["cargo", "run", "--locked", "--example", "yohaku_compositor"], project, sys.stdout.buffer, env)
    except (OSError, ValueError) as error:
        print(f"Could not launch Yohaku: {error}", file=sys.stderr)
        return 1


if __name__ == "__main__":
    raise SystemExit(main())
```

## 输出与关联
```yaml
output: "agent/yohaku_launcher.py"
```
